% Actual mixed Ferrari quotients; no finite-to-GT lifting assertion.
\subsection{An exact mixed boundary and compact quotient}

\paragraph{MF.1 (Proved): odd relative completions over a forgetful
base retain no Ferrari kernel.}
Use the fold $q:L=B_4/\langle z_4\rangle\to B$ and the
free kernel generators $\xi_1,\xi_2$ of FP.26. Put
$d=abc$ and $j=d^2$. In $L$, $j^2=1$, and
\[
 d\xi_1d^{-1}=\xi_2^{-1},\qquad
 d\xi_2d^{-1}=\xi_1.
\]
These follow by composing the two Nielsen actions in
FP.26.2 and using $c=\xi_1a$. Hence $j$ inverts both
free generators. Since $q(d^2)=z_3$ before quotienting the
three-strand centre, the kernel $\mathcal C=\ker(L\to B)$
is $F(\xi_1,\xi_2)\rtimes\langle j\rangle$.
With
\[
       t_0=j,\qquad t_1=\xi_1j,\qquad t_2=\xi_2j,
\]
its presentation is exactly
\[
               \mathcal C=C_2*C_2*C_2.
\]
The same statement holds for the indicated profinite
completions, by the split finite-extension description.

Let $N_L$ be the inverse image of an open normal
$N_B\triangleleft B$. For every odd prime $p$, a finite
$p$-group quotient of $N_L$ kills each $t_i$, since they
have order two. It therefore kills $\mathcal C$.
The universal property gives the exact isomorphisms
\[
 L/O^p(N_L)\simeq B/O^p(N_B),\qquad
 L/(N_L'N_L^p)\simeq B/(N_B'N_B^p).
\]
In the second equality one can also use $t_i^p=t_i$.
Thus even the full relative pro-$p$ completion over a
forgetful base, for odd $p$, cannot add a Ferrari coordinate.

\paragraph{MF.2 (Proved): the first mixed homology stage and its
precise pentagon limitation.}
Put $V=\mathcal C^{\rm ab}\simeq(\mathbb F_2)^3$ with
basis the three $t_i$. Its commutator subgroup is free of
rank five: the eight-sheeted Cayley covering is the cube,
with twelve unoriented edges and eight vertices.
Over $\mathbb F_3$, its homology is
\[
 H=\mathcal C'^{\rm ab}/3
   =B_0\oplus P_0,
 \qquad\dim B_0=3,\quad\dim P_0=2,
\]
with respective $V$ character masks
$(3,5,6)$ and $(7,7)$. A mask acts by minus one precisely
when its dot product with the base element is one.
Equivalently the three distinct pair-sign characters occur
once, and the all-minus character occurs twice.
The absence of other characters follows from the cellular
character formula
$2\mathbb F_3[V]-\sum_i\mathbb F_3[V/\langle t_i\rangle]
 +\mathbb F_3$.

An explicit model of
$\mathcal C/(\mathcal C''\mathcal C'^3)$ is
$H\rtimes V$, of order $8\cdot3^5=1944$, generated by
\[
 t_0=(0,e_0),\quad
 t_1=((1,0,0,1,0),e_1),\quad
 t_2=((0,1,1,0,1),e_2).
\]
The exact conjugation actions induced by $a,b$ are
\[
 \begin{split}
 a:(t_0,t_1,t_2)&\longmapsto(t_0t_1t_0,t_0,t_2),\\
 b:(t_0,t_1,t_2)&\longmapsto(t_2,t_1,t_2t_0t_2),
 \end{split}
\]
and $c$ acts as $\operatorname{Inn}(t_1t_0)$ followed by
the action of $a$. For example
$j=\xi_1^{-1}\xi_2^{-1}z_3$, so the Nielsen formulas give
$aja^{-1}=\xi_1^{-1}j$ and $bjb^{-1}=\xi_2j$.
This derives the displayed actions without a choice of
extension of an outer automorphism.

The affine action on $H$ splits into its two displayed
blocks. Each of the five pure $X_i$ has identity linear
part on $B_0$, so those five images are translations
there. Every evaluation of $f\in F'$ is therefore one on
that block. On $P_0$, simultaneous conjugation by the
translation $(1,1)$ gives exactly
\[
           a\mapsto(0,U),\qquad b\mapsto(0,L_0),
                    \qquad c\mapsto(e_1,U),
\]
the already audited order-$216$ affine quotient.
Adjoining $q_{12}$ through the fold does not alter this
pentagon test: its five $X_i$ images are
$(x,y,x,1,1)$, and the resulting pentagon is
$f(y,x)f(x,y)=1$ by antisymmetry.
In particular LI.8's actual witness passes this joint
test with $\Delta=1$. The exact affine matrices and all
source relators are checked in
\texttt{check\_mixed\_ferrari\_quotient.py}.

\paragraph{MF.3 (Proved construction): an actual class-two quotient
retaining the six mixed commutators.}
Let
\[
 Z=B_0\otimes P_0,\qquad\dim Z=6,
 \qquad N=H\oplus Z.
\]
Give $N$ the class-two exponent-three group law
\[
 (v,z)(v',z')=
 \left(v+v',\ z+z'+\frac12[v,v']\right),
 \quad
 [(b,p),(b',p')]=b\otimes p'-b'\otimes p.
\]
Here $1/2=2$ in $\mathbb F_3$.
The $V$ action on $Z$ has masks $(4,4,2,2,1,1)$.
Set $E=N\rtimes V$, of order
\[
                         |E|=8\cdot3^{11}=1\,417\,176.
\]
Use the same three $t_i$ as above, with zero $Z$
coordinates. They are involutions, since their degree-one
coordinates lie in the minus eigenspaces of the respective
base elements. They generate $E$: their commutator subgroup
surjects onto $H$, as checked in MF.2, and the Frattini
criterion for the finite three-group $N$ then gives
surjectivity onto $N$; equivalently their mixed brackets
span all of $Z$.

This is a canonical $L$-stable quotient of the full
class-two exponent-three quotient of $\mathcal C'$.
Indeed that full quotient has centre $\bigwedge^2H$.
The discarded subspace
$\bigwedge^2B_0\oplus\bigwedge^2P_0$ is stable under
the exact $a,b,c$ actions, because their degree-one
matrices preserve $B_0,P_0$ and the action on commutators
is their exterior square. The quotient retains exactly
the $B_0\otimes P_0$ component. Since the first-stage
group $H\rtimes V$ is centreless and $Z^V=0$, the group
$E$ is centreless. Its conjugation action therefore gives
an actual finite quotient of $L$, faithful on this
specified quotient of $\mathcal C$.

Here are complete small coordinates for that action.
Write $M_a=\operatorname{diag}(B_a,U)$,
$M_b=\operatorname{diag}(B_b,L_0)$,
$M_c=\operatorname{diag}(B_c,U)$ over $\mathbb F_3$, where
\[
 B_a=\begin{pmatrix}1&0&0\\0&0&1\\0&1&0\end{pmatrix},
 \quad
 B_b=\begin{pmatrix}0&0&2\\0&1&0\\2&0&0\end{pmatrix},
 \quad
 B_c=\begin{pmatrix}1&0&0\\0&0&2\\0&2&0\end{pmatrix}.
\]
The corresponding inner translations $n_g=(u_g,z_g)$
are
\[
\begin{array}{c|c|c}
g&u_g&z_g\\ \hline
a&(1,0,0,1,0)&(2,1,0,0,0,0)\\
b&(2,2,2,0,2)&(0,1,2,2,0,2)\\
c&(2,0,0,2,0)&(2,1,0,0,0,0).
\end{array}
\]
Thus $g$ acts as $\operatorname{Inn}(n_g)$ followed by
the graded action $M_g$ on $H$, $B_g\otimes P_g$ on
$Z$, and its induced permutation on $V$.
Equivalently these are elements $(n_g,M_g)$ of the finite
semidirect group with multiplication
\[
       (n,M)(n',M')=(n\,M(n'),MM').
\]
Their central coordinates are not guessed from an exterior
power representation. They are obtained by evaluating
the three exact involution images above with the BCH law,
then solving their central-coordinate equations. Each
central character is nontrivial, so the solution is unique;
the solution is checked on all three generators.

The code \texttt{check\_mixed\_ferrari\_class2.py} verifies
the four defining relations of $L$, as well as
\[
 (abc)^2=\operatorname{Inn}(t_0),\qquad
 ca^{-1}=\operatorname{Inn}(\xi_1),\qquad
 aba c^{-1}b^{-1}a^{-1}=\operatorname{Inn}(\xi_2)
\]
inside this exact model. It supplies all five $X_i$
coordinates and reusable group operations.
This paragraph constructs the genuine finite target.
It makes no assertion that its pentagon forces $\Delta=0$,
that a finite pentagon assignment lifts, or that its pair
second derived groups are nonnilpotent; those are separate
questions about this target.

\paragraph{MF.4 (Proved): one actual mixed-sign element satisfies the
pentagon in the full relative pronilpotent Ferrari quotient.}
Write
\[
 C_\ell=\widehat{\mathcal C}/O^\ell(\widehat{\mathcal C}'),
 \qquad
 L_\ell=\widehat L/O^\ell(\widehat{\mathcal C}').
\]
There exists one actual $k=(1,f_k)\in\mathcal H$, retaining
all properties of LI.8 and LI.9, such that
\[
 \boxed{\quad\Delta(k)=1,\qquad
 \mathcal P(f_k)=1\ \text{in}\
       \widehat L/R_{\rm nil}(\widehat{\mathcal C}').\quad}
\]
Thus the pentagon is satisfied simultaneously in every
$L_\ell$, with its entire four-point quotient retained.
This does not identify the relative pronilpotent quotient
with the full five-point group, or assert the full pentagon.

\emph{Exact integral linear input.}
The cube-graph construction in MF.2 works over $\mathbb Z$.
Conjugation by $\xi_1,\xi_2$ on its rank-five cycle lattice
gives commuting involutions $T_1,T_2$. Put
\[
 E=\frac{1+T_1+T_2+T_1T_2}{4}.
\]
This is a rank-two projector over $\mathbb Z[1/2]$,
commuting with the $a,b,c$ actions; it is the compact
summand. If $X_i$ denotes the corresponding integral
homology action, direct integral identities give
\[
 \begin{split}
 (X_i-I)(I-E)&=0&& (1\leq i\leq5),\\
 (X_i-I)^2E&=0&& (i=1,2,3),\\
 (X_i+I)E&=0&& (i=4,5).
 \end{split}
\]
These are verified over the integers after multiplying
by four, not inferred from a computation modulo three.
The source is the cube with eight vertices, twelve edges
and five non-tree edges; its five explicit free loops and
the exact Nielsen substitutions are recorded by
\texttt{check\_integral\_ferrari\_homology.py}.
Consequently, at every odd prime $\ell$, the first three
linear $X_i$ have order dividing $\ell$ and the last two
have order dividing two.

For odd $\ell$, the free pro-$\ell$ normal subgroup of
$C_\ell$ has rank five and quotient $V=(C_2)^3$. The
$V$ action on its Frattini quotient is faithful by the
five character masks of MF.2. Its centralizer in $C_\ell$
is therefore trivial: a centralizing element has trivial
$V$ coordinate, and then lies in the centre of a
nonabelian free pro-$\ell$ group. Restriction consequently
embeds $\operatorname{Aut}(C_\ell)$ in the automorphism
group of that free pro-$\ell$ subgroup. The kernel of
its action on the Frattini quotient is pro-$\ell$.
This is the usual finite $p$-group Frattini-automorphism
theorem, passed to the inverse limit of characteristic
finite quotients. Thus the procyclic closure of each
$X_i$ in this automorphism group has primes only in
$\{2,\ell\}$.

At $\ell=2$, $C_2$ is the pro-two completion of
$\mathcal C$. The finite-index two-group extension
description identifies it with
$\widehat F_2^{(2)}\rtimes C_2$, with the involution
inverting the two free generators. It is centreless:
the nontrivial coset acts as $-I$ on
$\mathbb Z_2^2$, whereas inner automorphisms act trivially
there, and the free pro-two group has trivial centre.
The action on its three-dimensional Frattini quotient
permutes the three $t_i$; every pure $X_i$ acts trivially.
The same Frattini-automorphism theorem therefore makes
each $X_i$ pro-two.

\emph{The exceptional prime three.}
Here the whole-family argument has an additional exact
central ingredient. In the closure $G$ of $\langle a,b\rangle$
inside $\operatorname{Aut}(C_3)$, the linear image is
$S_3\times\operatorname{SL}_2(\mathbb F_3)$, whose
second derived subgroup is the central $C_2$ in the
second factor. The actual central braid $z_3=(ab)^3$
maps to its generator. Since the linear kernel is
pro-three, the two-primary component of $z_3$ is a
central involution, and the full preimage of that
$C_2$ is a direct product of a pro-three group and
$C_2$. It is pronilpotent. The second derived group
of the harmonic pair lies in this preimage; the other
harmonic pair is conjugate by the reversing half twist.
For the remaining three pairs, the linear images are
abelian, so their derived groups are pro-three.
It follows that any $f\in R_{\rm nil}(F'')$ has all
five factor evaluations equal to one in
$\operatorname{Aut}(C_3)$.

\emph{A single uniform actual lifting seed.}
Let $e_3\in\widehat{\mathbb Z}$ have component one in
$\mathbb Z_3$ and zero at every other prime. Set
\[
        a_*=[[x^{e_3},y^{e_3}],[y^{e_3},z^{e_3}]].
\]
This is a genuine profinite word in $F''$.
Since $e_3\equiv4\pmod{12}$, its original $q_{12}$
value is $[5I]$. Every full orbifold conjugate kills
it because its order-two peripheral has $e_3$th power
one. For every $\ell\ne3$, the preceding procyclic
order statements give $X_i^{e_3}=1$ in
$\operatorname{Aut}(C_\ell)$. Therefore $a_*$ vanishes
at every one of the five pair evaluations and all
six peripheral permutations; hence it belongs to
all their $B$-normal core kernels.

For a finite list of finite quotients of these
automorphism targets, intersect those core kernels
with the $N$ used in LI.8. The same seed $a_*$ belongs
to every such $N$ and retains $q_{12}(a_*)=[5I]$.
The coprime centralizer lift and actual R2 lifting
therefore give the required finite intersection
property, including the normalization conditions
$f,\beta\in R_{\rm nil}(N)$. Compactness yields an
actual $k$ retaining the LI.8 frames, full orbifold
kernel and $R_{\rm nil}(F'')$ condition, with all
five evaluations one in every
$\operatorname{Aut}(C_\ell)$ for $\ell\ne3$.
The preceding three-adic argument supplies the
remaining prime. The earlier LI.9 conditions may be
included simultaneously: $a_*$ kills its three
nonharmonic evaluation kernels away from three,
and its harmonic and three-adic conditions follow
from $R_{\rm nil}(F'')$ just as before.

\emph{From the action back to the actual source quotient.}
For any multiplier-one antisymmetric pair, forgetting
sends $(X_1,X_2,X_3,X_4,X_5)$ to $(x,y,x,1,1)$,
so its pentagon becomes $f(y,x)f(x,y)=1$ in
$\widehat B$. Thus the pentagon of the constructed
element belongs to $C_\ell$ inside $L_\ell$.
All five factor evaluations act trivially on $C_\ell$,
so their product centralizes it. Since $C_\ell$ is
centreless, that pentagon is one in $L_\ell$.
This step asserts factorwise vanishing in the
automorphism group and pentagon vanishing in the
actual source quotient; it does not confuse them.
Finally
\[
 \bigcap_\ell O^\ell(\widehat{\mathcal C}')
             =R_{\rm nil}(\widehat{\mathcal C}')
\]
by the Sylow-product description of finite nilpotent
groups. This proves the boxed simultaneous statement.
